\documentclass[12pt]{article}

\usepackage[margin=1in]{geometry}
\usepackage{amsmath,amssymb,mathtools}
\usepackage{enumitem}
\usepackage{comment}
\usepackage{hyperref}
\usepackage{xcolor}

% -----------------------
%  Show / hide solutions
% -----------------------
\newif\ifshowAns
% \showAnstrue      % <-- instructor version
\showAnsfalse   % <-- student version

\ifshowAns
  \newenvironment{sol}{\par\smallskip\noindent \textcolor{red}{\textbf{Solution.}} \\ }{\par\medskip}
\else
  \excludecomment{sol}
\fi

\newcommand{\pts}[1]{\hfill{\small[#1 pts]}}
\newcommand{\E}{\mathbb{E}}
\newcommand{\R}{\mathbb{R}}
\newcommand{\norminf}[1]{\left\lVert #1 \right\rVert_{\infty}}

\begin{document}

\begin{center}
{\Large \textbf{Midterm Exam 1}}\\
Seminar in Macroeconomics\\
\medskip
\textit{Hui-Jun Chen}
\end{center}

% ============================================================
\section*{Question 1 \pts{35}}

\subsection*{Environment}
Consider the deterministic neoclassical growth problem with no labor:
\[
\max_{\{c_t,k_{t+1}\}_{t\ge0}} \sum_{t=0}^{\infty}\beta^t \log(c_t),\qquad \beta\in(0,1)
\]
subject to
\[
c_t + k_{t+1} = k_t^\alpha + (1-\delta)k_t,\qquad c_t\ge 0,\ k_{t+1}\ge 0,
\]
with given \(k_0>0\), \(\alpha\in(0,1)\), and \(\delta\in(0,1)\).

\begin{enumerate}[label=(\alph*)]
\item \textbf{Three representations (conceptual).} \pts{12}\\
In \textbf{one short paragraph each}, describe what changes across:
\begin{itemize}
\item Date-0 (plan) representation,
\item Sequential representation,
\item Recursive (DP) representation.
\end{itemize}
You do \emph{not} need to write the full math program, but you must clearly say what the \textbf{choice object} is in each case.

\item \textbf{Write the Bellman equation.} \pts{10}\\
Define the state and choice and write the Bellman equation \(V(k)=\max\{\cdots\}\).

\item \textbf{Why recursion is computationally useful.} \pts{8}\\
Explain (in words) why the recursive formulation can be solved by iteration \(V_{n+1}=TV_n\),
while the Date-0 and Sequential formulations are awkward to put directly on a computer.

\item \textbf{Blackwell intuition (no proof).} \pts{5}\\
In one or two sentences, explain why discounting by \(\beta<1\) suggests that \(T\) ``shrinks'' differences
between value functions (geometric intuition).
\end{enumerate}

\begin{sol}
(a) \textbf{Date-0:} choose entire sequences \(\{c_t,k_{t+1}\}_{t\ge0}\) at time 0 subject to all feasibility constraints.
\textbf{Sequential:} still choose sequences, but constraints are written period-by-period (plus a TVC/no-Ponzi condition).
\textbf{Recursive:} choose a \emph{rule} \(k'=g(k)\) state-by-state; problem is summarized by a Bellman equation where the
unknown is the function \(V(\cdot)\).

(b) State: \(k\). Choice: \(k'\in[0,\,k^\alpha+(1-\delta)k]\). Consumption: \(c=k^\alpha+(1-\delta)k-k'\).
\[
V(k)=\max_{k'\in[0,\,k^\alpha+(1-\delta)k]}
\left\{\log\!\big(k^\alpha+(1-\delta)k-k'\big)+\beta V(k')\right\}.
\]

(c) DP turns an infinite-horizon plan into a fixed point \(V=TV\) over functions; we can iterate the same one-period problem.
Date-0 involves infinite sequences and infinite sums directly; sequential involves infinitely many constraints/multipliers.
Recursive reduces the problem to repeatedly applying the same operator until convergence.

(d) If two guesses differ by at most \(m\), then next period continuation differs by at most \(\beta m\), then \(\beta^2 m\), etc.
Geometric discounting makes distant differences shrink quickly, supporting contraction/convergence.
\end{sol}

% ============================================================
\section*{Question 2: One-sided Labor Search \pts{35}}

An unemployed worker receives an i.i.d.\ wage offer \(w\sim F\) each period on \([0,B]\).
If she \textbf{accepts} an offer \(w\), she earns wage \(w\) forever.
If she \textbf{rejects}, she receives unemployment payoff \(c>0\) for the period and draws a new offer next period.
Preferences are risk-neutral: \(\sum_{t\ge0}\beta^t c_t\).

\begin{enumerate}[label=(\alph*)]
\item \textbf{Bellman equation and accept/reject values.} \pts{12}\\
Write \(v(w)\) and show the accept value equals \(\frac{w}{1-\beta}\).
Write the reject value in terms of \(\int v(w')\,dF\).

\item \textbf{Reservation wage rule.} \pts{10}\\
Define the reservation wage \(\bar w\) and state the accept/reject rule in terms of \(\bar w\).
(You may give an indifference condition.)

\item \textbf{Comparative statics (intuition).} \pts{8}\\
Explain why \(\bar w\) increases with \(c\) and with \(\beta\). (Words are fine.)

\item \textbf{Duration.} \pts{5}\\
Let \(H=1-F(\bar w)\) be the probability of accepting in a given period.
Under i.i.d.\ offers, what is the expected unemployment duration?
\end{enumerate}

\begin{sol}
(a) Bellman:
\[
v(w)=\max\left\{\underbrace{w+\beta w+\beta^2 w+\cdots}_{=w/(1-\beta)},\;
\underbrace{c+\beta\int_0^B v(w')\,F(dw')}_{\text{reject}}\right\}.
\]
So accept value is \(w/(1-\beta)\), reject value is \(c+\beta\E[v(w')]\).

(b) \(\bar w\) solves indifference:
\[
\frac{\bar w}{1-\beta}=c+\beta\int v(w')\,F(dw').
\]
Rule: accept iff \(w\ge \bar w\), reject iff \(w<\bar w\).

(c) Higher \(c\) improves the outside option of rejecting, so you become pickier \(\Rightarrow \bar w\uparrow\).
Higher \(\beta\) raises the value of waiting for a better offer, increasing the option value of rejecting \(\Rightarrow \bar w\uparrow\).

(d) Each period you accept with probability \(H\). Duration is geometric with mean \(1/H\).
\end{sol}

% ============================================================
\section*{Question 3: Two-sided Labor Search \pts{30}}

Consider a Mortensen--Pissarides labor market.
Unemployment is \(u\), vacancies are \(v\), and matching is \(M(u,v)\) with CRS.
Define tightness \(\theta=v/u\).
Let the vacancy-filling probability be \(q(\theta)=M(u,v)/v\) and job-finding probability be \(f(\theta)=M(u,v)/u=\theta q(\theta)\).
Suppose output per match is \(z\), separation rate is \(\delta\), vacancy posting cost is \(p\), and unemployment payoff is \(c\).
Wage is determined by Nash bargaining with worker bargaining power \(\phi\in(0,1)\).

\begin{enumerate}[label=(\alph*)]
\item \textbf{Matching rates.} \pts{8}\\
Define \(q(\theta)\) and \(f(\theta)\). Explain (in words) why typically \(q'(\theta)<0\) and \(f'(\theta)>0\).

\item \textbf{Free entry condition.} \pts{10}\\
Let \(J_u\) be the value of an unmatched firm (vacancy) and \(J_e\) be the value of a matched firm.
State the free entry condition and show it implies a relationship of the form
\[
p=\beta q(\theta)J_e.
\]
(You do not need to derive the full Bellman equations; a short argument is fine.)

\item \textbf{Wage equation.} \pts{8}\\
State the Nash bargaining implication that wage can be written as
\[
w(\theta)=\phi z + (1-\phi)c + \phi\theta p.
\]
Explain in one sentence why \(w\) increases with \(\theta\).

\item \textbf{Steady-state unemployment.} \pts{4}\\
Write the steady-state unemployment rate in terms of \(\delta\) and \(f(\theta)\).
\end{enumerate}

\begin{sol}
(a) \(q(\theta)=M(u,v)/v\) and \(f(\theta)=M(u,v)/u=\theta q(\theta)\).
With more vacancies per unemployed worker (\(\theta\uparrow\)), it becomes harder for each vacancy to find a worker
(\(q(\theta)\downarrow\)), but easier for each unemployed worker to find a job (\(f(\theta)\uparrow\)).

(b) Free entry means firms post vacancies until the value of an open vacancy is zero: \(J_u=0\).
From the vacancy value equation \(J_u=-p+\beta[q(\theta)J_e+(1-q(\theta))J_u]\),
setting \(J_u=0\) yields \(0=-p+\beta q(\theta)J_e\Rightarrow p=\beta q(\theta)J_e\).

(c) Nash bargaining splits match surplus; rearranging gives
\(w(\theta)=\phi z+(1-\phi)c+\phi\theta p\).
Higher \(\theta\) means workers find jobs more easily and firms fill vacancies more slowly, raising the worker’s outside option/surplus share,
so wages rise in tight markets.

(d) Steady state satisfies flows in = flows out:
\(\delta(1-u)=uf(\theta)\Rightarrow u=\frac{\delta}{\delta+f(\theta)}\).
\end{sol}

\end{document}
